\documentclass[12pt]{article}
\usepackage{amsmath, amssymb}

\begin{document}

\section*{Lecture Scribe}

\textbf{CSE400 – Fundamentals of Probability in Computing} \\
\textbf{Lecture 18: Marginal PMF, Correlation, Covariance, Joint Gaussian Random Variables} \\
\textbf{Instructor: Dhaval Patel, PhD} \\
\textbf{Date: March 12, 2026}

\section{Marginal PMF}

\subsection{Definition}

For discrete random variables $X$ and $Y$, the joint PMF is defined as:
\[
p_{XY}(x,y) = P(X = x, Y = y)
\]

The marginal PMFs are obtained by summing over the other variable:
\[
p_X(x) = \sum_y p_{XY}(x,y)
\]
\[
p_Y(y) = \sum_x p_{XY}(x,y)
\]

\textbf{Key Idea:} The marginal PMF of one variable is obtained by summing the joint PMF over all possible values of the other variable.

\subsection{Marginal PMF from a Joint PMF Table}

From the table representation (page 4):

\begin{itemize}
    \item Row sums give $p_X(x)$
    \item Column sums give $p_Y(y)$
\end{itemize}

Mathematically:
\[
p_X(x_i) = \sum_j p_{ij}, \quad p_Y(y_j) = \sum_i p_{ij}
\]

Also:
\[
\sum_i p_X(x_i) = 1, \quad \sum_j p_Y(y_j) = 1
\]

\subsection{Solved Example}

Given joint PMF (page 5):

\[
\begin{array}{c|cc}
 & y=0 & y=1 \\
\hline
x=0 & \frac{1}{8} & \frac{1}{4} \\
x=1 & \frac{3}{8} & \frac{1}{4}
\end{array}
\]

\textbf{Step 1: Compute $p_X(x)$}

\[
p_X(0) = \frac{1}{8} + \frac{1}{4} = \frac{3}{8}
\]
\[
p_X(1) = \frac{3}{8} + \frac{1}{4} = \frac{5}{8}
\]

\textbf{Step 2: Compute $p_Y(y)$}

\[
p_Y(0) = \frac{1}{8} + \frac{3}{8} = \frac{1}{2}
\]
\[
p_Y(1) = \frac{1}{4} + \frac{1}{4} = \frac{1}{2}
\]

Hence:
\[
p_X(x) =
\begin{cases}
\frac{3}{8}, & x = 0 \\
\frac{5}{8}, & x = 1
\end{cases}
\quad
p_Y(y) =
\begin{cases}
\frac{1}{2}, & y = 0 \\
\frac{1}{2}, & y = 1
\end{cases}
\]

\section{Correlation}

\subsection{Definition}

A basic way to measure joint behavior is through the product moment:
\[
R_{XY} = E[XY]
\]

It represents the average value of the product $XY$.

It is computed about the origin (not about the means).

\subsection{Interpretation}

Depends on the scale of $X$ and $Y$.

Mixes:
\begin{itemize}
    \item mean values
    \item spread
    \item dependence
\end{itemize}

If:
\[
R_{XY} = 0
\]
then variables are orthogonal about the origin.

\textbf{Limitation:} It is not a clean measure of centered dependence.

\textbf{Transition:} To remove the effect of means, variables are centered $\rightarrow$ leads to covariance.

\section{Covariance}

\subsection{Definition and Interpretation}

Covariance is defined as:
\[
\text{Cov}(X,Y) = E[(X - \mu_X)(Y - \mu_Y)]
\]

where:
\[
\mu_X = E[X], \quad \mu_Y = E[Y]
\]

\textbf{Interpretation:}
\begin{itemize}
    \item Measures how centered variables vary together
    \item Indicates direction of joint linear variation:
    \begin{itemize}
        \item $> 0$: positive covariance
        \item $< 0$: negative covariance
        \item $= 0$: no linear co-variation
    \end{itemize}
\end{itemize}

\textbf{Important Note:} Magnitude depends on units and cannot measure strength.

\subsection{Algebraic Form}

Starting with:
\[
\text{Cov}(X,Y) = E[(X - \mu_X)(Y - \mu_Y)]
\]

Expanding:
\[
= E[XY - X\mu_Y - \mu_X Y + \mu_X \mu_Y]
\]
\[
= E[XY] - \mu_Y E[X] - \mu_X E[Y] + \mu_X \mu_Y
\]

Since:
\[
\mu_X = E[X], \quad \mu_Y = E[Y]
\]

We obtain:
\[
\text{Cov}(X,Y) = E[XY] - E[X]E[Y]
\]

\subsection{Properties of Covariance}

\begin{itemize}
    \item \textbf{Symmetry}
    \[
    \text{Cov}(X,Y) = \text{Cov}(Y,X)
    \]
    
    \item \textbf{Variance as Special Case}
    \[
    \text{Cov}(X,X) = \text{Var}(X)
    \]
    
    \item \textbf{Linearity}
    \[
    \text{Cov}(aX + b, cY + d) = ac \, \text{Cov}(X,Y)
    \]
    
    \item \textbf{Independence}
    \[
    X \perp Y \Rightarrow \text{Cov}(X,Y) = 0
    \]
\end{itemize}

\textbf{Important:} Zero covariance $\Rightarrow$ uncorrelated, but does not imply independence.

\subsection{Example 1}

Given:
\[
Y = -2X + 3
\]

Using:
\[
\text{Cov}(X,Y) = E[XY] - E[X]E[Y]
\]

Compute:
\[
XY = X(-2X + 3) = -2X^2 + 3X
\]
\[
E[Y] = E[-2X + 3] = -2E[X] + 3
\]

Substituting:
\[
\text{Cov}(X,Y) = E[-2X^2 + 3X] - E[X](-2E[X] + 3)
\]

Simplifying:
\[
= -2E[X^2] + 3E[X] + 2(E[X])^2 - 3E[X]
\]
\[
= -2E[X^2] + 2(E[X])^2
\]

Using:
\[
\text{Var}(X) = E[X^2] - (E[X])^2
\]

We obtain:
\[
\text{Cov}(X,Y) = -2 \, \text{Var}(X)
\]

\subsection{Pearson’s Correlation Coefficient}

\[
\rho_{XY} = \frac{\text{Cov}(X,Y)}{\sqrt{\text{Var}(X)\text{Var}(Y)}} = \frac{\text{Cov}(X,Y)}{\sigma_X \sigma_Y}
\]

\textbf{Properties:}
\begin{itemize}
    \item Measures strength and direction of linear relationship
    \item Normalized version of covariance
    \item Unit-free
    \item Range:
    \[
    -1 \le \rho_{XY} \le 1
    \]
\end{itemize}

\textbf{Interpretation:}
\begin{itemize}
    \item $\rho_{XY} > 0$: positive linear trend
    \item $\rho_{XY} < 0$: negative linear trend
    \item $\rho_{XY} = 0$: no linear correlation
\end{itemize}

\section{Example 2 (Continuous Case)}

Given joint PDF (page 14):
\[
f_{XY}(x,y) =
\begin{cases}
\frac{xy}{96}, & 0 < x < 4, \; 1 < y < 5 \\
0, & \text{otherwise}
\end{cases}
\]

\textbf{Results from Solution}

\[
E[X] = \frac{8}{3}
\]
\[
E[Y] = \frac{31}{9}
\]

Since $X$ and $Y$ are independent:
\[
E[XY] = E[X]E[Y] = \frac{248}{27}
\]

\[
E[2X + 3Y] = \frac{47}{3}
\]

\[
\text{Var}(X) = \frac{8}{9}
\]
\[
\text{Var}(Y) = \frac{92}{81}
\]

\[
\text{Cov}(X,Y) = 0
\]
\[
\rho_{XY} = 0
\]

\section{Jointly Gaussian Random Variables}

\subsection{Definition}

A pair $(X,Y)$ is jointly Gaussian if:
\begin{itemize}
    \item The joint density has a bivariate Gaussian form
    \item The marginals are Gaussian
\end{itemize}

Applications:
\begin{itemize}
    \item Signal processing
    \item Communication systems
    \item Modeling continuous noisy measurements
\end{itemize}

\subsection{Joint Gaussian PDF}

\[
f_{X,Y}(x,y) =
\frac{1}{2\pi \sigma_X \sigma_Y \sqrt{1 - \rho^2}}
\exp\left(
-\frac{1}{2(1 - \rho^2)}
\left[
\left(\frac{x - \mu_X}{\sigma_X}\right)^2
+
\left(\frac{y - \mu_Y}{\sigma_Y}\right)^2
-
2\rho
\left(\frac{x - \mu_X}{\sigma_X}\right)
\left(\frac{y - \mu_Y}{\sigma_Y}\right)
\right]
\right)
\]

where:
\begin{itemize}
    \item $\mu_X, \mu_Y$: means
    \item $\sigma_X, \sigma_Y$: standard deviations
    \item $\rho$: correlation coefficient
\end{itemize}

\subsection{Interpretation}

Example (page 23):

\begin{itemize}
    \item $X$: Atmospheric CO$_2$ concentration
    \item $Y$: Global surface temperature anomaly
\end{itemize}

Because both variables are influenced by multiple random processes, their joint behavior can be approximated by a joint Gaussian distribution.

If:
\[
\rho > 0
\]
then higher CO$_2$ levels tend to occur with higher temperatures.

\section*{End of Lecture 18 Scribe}

\end{document}